\chapter{Worked example: ElGamal encryption}

Following Rosulek, \emph{The Joy of Cryptography} (ElGamal / DDH). ElGamal is the
first \emph{computational} example: over an abstract prime-order group its IND-CPA
security is not perfect but reduces to the decisional Diffie--Hellman assumption,
with the standard factor-two bound. This exercises the reduction machinery rather
than a bijection coupling.

\begin{definition}[ElGamal over a prime-order group]
  \label{def:elgamal}
  \uses{def:spcomp}
  \lean{CatCrypt.Examples.ElGamal.elgamalEnc, CatCrypt.Examples.ElGamal.elgamalDDH}
  \leanok
  With public key $\mathit{pk} = g^{x}$, ElGamal encrypts $m$ as $(g^{r}, m \cdot
  \mathit{pk}^{r})$ for fresh $r$. The matching DDH instance fixes
  $\mathsf{ecdh}(g^{a}, g^{b}) = g^{ab}$ on group elements.
\end{definition}

\begin{theorem}[Correctness]
  \label{thm:elgamal-correct}
  \uses{def:elgamal}
  \lean{CatCrypt.Examples.ElGamal.elgamal_dec_enc}
  \leanok
  Decryption inverts encryption: $(m \cdot \mathit{pk}^{r}) \cdot (g^{r})^{-x} =
  m$.
\end{theorem}

\begin{theorem}[IND-CPA reduces to DDH]
  \label{thm:elgamal-ddh}
  \uses{def:elgamal}
  \lean{CatCrypt.Examples.ElGamal.elgamal_indcpa_le_ddh}
  \leanok
  ElGamal's IND-CPA advantage is bounded by the DDH advantage of two explicit
  reduction adversaries:
  \[
    \mathsf{Adv}^{\mathrm{IND\text{-}CPA}}_{\mathrm{ElGamal}}(m_0, m_1, A) \;\le\;
    \mathsf{Adv}^{\mathrm{DDH}}(R_{m_0}) + \mathsf{Adv}^{\mathrm{DDH}}(R_{m_1}).
  \]
  The three-hop argument passes through the game where the mask is uniform, in
  which the message is absorbed and the two challenges coincide.
\end{theorem}

\begin{theorem}[Exact real-or-random bound]
  \label{thm:elgamal-eq-ddh}
  \uses{def:elgamal}
  \lean{CatCrypt.Examples.ElGamal.elgamal_indcpa_eq_ddh}
  \leanok
  In the real-or-random framing (real encryption versus a random ciphertext), the
  advantage \emph{equals} a single DDH advantage --- no factor and no sum:
  \[
    \mathsf{Adv}(\mathrm{Enc}\,m_0,\ \mathrm{random}) \;=\;
    \mathsf{Adv}^{\mathrm{DDH}}(R_{m_0}).
  \]
  This matches SSProve-Rocq's exact one-time-secrecy result; the two-message
  left-or-right advantage above is a different, two-term quantity.
\end{theorem}
